% Generated draft: requires mathematical and source review.
\documentclass[11pt]{article}
\usepackage[margin=25mm]{geometry}
\usepackage{fontspec}
\setmainfont{Latin Modern Roman}
\usepackage{amsmath,amssymb,mathtools,unicode-math}
\setmathfont{Latin Modern Math}
\usepackage{enumitem,needspace,xurl,hyperref}
\setlength{\parindent}{0pt}
\setlength{\parskip}{5pt}
\setlength{\emergencystretch}{3em}
\newcommand{\N}{\mathbb{N}}
\newcommand{\Z}{\mathbb{Z}}
\newcommand{\Q}{\mathbb{Q}}
\newcommand{\R}{\mathbb{R}}
\newcommand{\C}{\mathbb{C}}
\newcommand{\F}{\mathbb{F}}
\newcommand{\A}{\mathbb{A}}
\newcommand{\PP}{\mathbb{P}}
\newcommand{\E}{\mathbb{E}}
\newcommand{\eps}{\varepsilon}
\newcommand{\dd}{\,\mathrm{d}}
\begin{document}
\section*{UnsolvedMath ID: 30003276}
\textbf{Source problem code:} catalogue entry
\subsection{Definitions and mathematical statement}
The record asks whether the Brauer--Manin obstruction is the only obstruction to the Hasse principle for smooth Fano hypersurfaces. It does not specify the ground number field, the dimensions and degrees under consideration, or the family underlying its density formulation. The following gives the standard arithmetic formulation with the ground field left as an explicit parameter; its intended scope requires clarification.

Let \(k\) be a number field, let \(\Omega_k\) be its set of places, and let \(k_v\) be its completion at \(v\). Let \(X\subseteq\mathbb P^n_k\) be a smooth, projective, geometrically integral hypersurface of positive dimension. The Fano condition means that its anticanonical line bundle \(-K_X\) is ample. For a smooth hypersurface of degree \(d\), adjunction gives \(-K_X\cong\mathcal O_X(n+1-d)\), so this condition is equivalent to \(d\leq n\). The record gives no narrower restrictions on \(n\) or \(d\).

Write \(\mathbb A_k\) for the ring of adeles of \(k\). Since \(X\) is projective, its adelic points may be identified with \(X(\mathbb A_k)=\prod_{v\in\Omega_k}X(k_v)\). Define the cohomological Brauer group by
\[
\operatorname{Br}(X)=H^2_{\mathrm{\acute et}}(X,\mathbb G_m)_{\mathrm{tors}}.
\]
For \(\alpha\in\operatorname{Br}(X)\) and \(P_v\in X(k_v)\), evaluation gives \(\alpha(P_v)\in\operatorname{Br}(k_v)\). With \(\operatorname{inv}_v:\operatorname{Br}(k_v)\to\mathbb Q/\mathbb Z\) the local invariant map, define
\[
X(\mathbb A_k)^{\operatorname{Br}}=
\left\{(P_v)_v\in X(\mathbb A_k):
\sum_{v\in\Omega_k}\operatorname{inv}_v\bigl(\alpha(P_v)\bigr)=0
\text{ for every }\alpha\in\operatorname{Br}(X)\right\}.
\]
The sums have only finitely many nonzero terms. The Hasse principle for \(X\) is the implication \(X(\mathbb A_k)\ne\varnothing\Rightarrow X(k)\ne\varnothing\). A Brauer--Manin obstruction to this principle occurs when \(X(\mathbb A_k)\ne\varnothing\) but \(X(\mathbb A_k)^{\operatorname{Br}}=\varnothing\).

For a specified ground field \(k\) and a specified class of these hypersurfaces, the conjectural assertion attributed in the record to Colliot--Th\'el\`ene is
\[
\text{for every }X\text{ in that class},\qquad
X(\mathbb A_k)^{\operatorname{Br}}\ne\varnothing
\ \Longrightarrow\ X(k)\ne\varnothing.
\]
The supplied record does not determine whether \(k\) is intended to be \(\mathbb Q\), another fixed number field, or universally quantified over number fields.

The record also asks whether
\[
\delta_{\mathrm{global}}=\delta_{\mathrm{local}}.
\]
Neither density is defined in the supplied material: the parameter space, counting convention, height or ordering, normalization, and existence of the limits are unspecified. Moreover, the original statement describes this equality as a consequence, whereas the cleaned statement calls it equivalent. Brauer--Manin sufficiency alone does not imply equality between densities of rationally soluble and everywhere locally soluble hypersurfaces: locally soluble hypersurfaces with a Brauer--Manin obstruction must also be accounted for. Conversely, equality of densities can allow exceptional counterexamples to a pointwise assertion. Thus the claimed consequence or equivalence cannot be adopted without additional hypotheses and definitions.
\subsection{Short English statement}
The question attributed to Colliot--Th\'el\`ene asks whether every smooth Fano hypersurface with an adelic point orthogonal to its entire Brauer group has a rational point. The intended ground field and scope remain unspecified. The accompanying equality \(\delta_{\mathrm{global}}=\delta_{\mathrm{local}}\) requires missing density definitions and additional justification linking it to the pointwise conjecture.
\subsection{Sources}
\begin{itemize}
\item Supplied local UnsolvedMath registry record, permanent ID \(30003276\), title ``Brauer--Manin Sufficiency for Smooth Fano Hypersurfaces,'' exact source problem code \texttt{catalogue entry}, in the Oberwolfach Reports Open Problems collection. The record attributes the conjecture to Colliot--Th\'el\`ene; it supplies no named proposer for this entry.
\item Supplied report citation: \emph{Analytic Number Theory (2017)}, Oberwolfach Reports, pp.\ 2975--3030, DOI \texttt{10.4171/owr/2016/53}, \url{https://doi.org/10.4171/owr/2016/53}. This is also the supplied literature source labelled ``OWR 53 report.'' The record gives proposed year 2016; the relationship between that year and the citation's 2017 label has not been independently checked. No problem-specific page or original conjecture citation is supplied.
\item Supplied literature source: Browning, Le Boudec, and Sawin, \url{https://doi.org/10.4007/annals.2023.197.3.3}. The record supplies no title or precise theorem statement identifying its application here.
\item Supplied literature assessment dated exactly \texttt{2026-08-21}: the general Colliot--Th\'el\`ene conjecture invoked by the record remains unresolved despite strong almost-all results for random Fano hypersurfaces. This is supplied dated curation derived from the OWR material, not an independently verified literature assessment. The record itself requests independent verification of time-sensitive status and cited claims.
\item The registry labels the problem \texttt{open} and its statement \texttt{corrected\_verified}, with a supplied explanation that notation, typography, or omitted definitions were repaired while preserving mathematical content. These are source labels, not independent verification of the formulation, its density equivalence, or its current mathematical status.
\end{itemize}
\end{document}
